% LaTeX-header for faglige dokumenter (rapport, vedlegg, notater) — LOG650
% Brukes via: pandoc --include-in-header=dokument-header.tex

% --- Pakker ---
\usepackage{booktabs}        % pene tabell-streker
\usepackage{longtable}       % tabeller over flere sider
\usepackage{array}           % kolonne-spesifikasjoner
\usepackage{makecell}        % multi-linjet tabellinnhold
\usepackage{colortbl}        % farger i tabeller
\usepackage{xcolor}          % farger
\usepackage{graphicx}        % figurer
\usepackage{float}           % figurplassering
\usepackage{caption}         % figurtekst
\usepackage{fancyhdr}        % topp/bunntekst
\usepackage{lastpage}        % sidetall x av y
\usepackage{microtype}       % bedre typografi
% (babel lastes av pandoc via -V lang=nb-NO)
\usepackage{amsmath, amssymb} % matematikk

% --- Sidehoder/-foter ---
\pagestyle{fancy}
\fancyhf{}
\renewcommand{\headrulewidth}{0.3pt}
\renewcommand{\footrulewidth}{0.3pt}
\fancyhead[L]{\small\itshape LOG650 · Rune Grødem}
\fancyhead[R]{\small\itshape Vår 2026}
\fancyfoot[L]{\small rune.grodem@rogbr.no}
\fancyfoot[R]{\small Side \thepage{} av \pageref{LastPage}}

% --- Tabeller med streker i alle bokser ---
\renewcommand{\arraystretch}{1.25}      % luft i celler
\setlength{\tabcolsep}{6pt}              % horisontal padding

% Tving alle pandoc-tabeller til å bruke streker (override default booktabs-stil)
\let\oldtoprule\toprule
\let\oldmidrule\midrule
\let\oldbottomrule\bottomrule
\renewcommand{\toprule}{\hline}
\renewcommand{\midrule}{\hline}
\renewcommand{\bottomrule}{\hline}

% Vertikale streker mellom kolonner via redefinering av longtable-omgivelser
\AtBeginEnvironment{longtable}{\setlength{\arrayrulewidth}{0.4pt}}

% --- Figurer alltid sentrert med tekst under ---
% labelformat=empty: figurnummer skrives manuelt i caption-teksten (f.eks. "Figur 4.1: ..."),
% slik at LaTeX ikke prefikser med eget "Figur N." og dobler nummereringen.
\captionsetup{font=small,labelfont=bf,labelsep=none,labelformat=empty,justification=centering}
\setkeys{Gin}{width=0.85\linewidth,keepaspectratio}

% --- Lenker ---
\usepackage{hyperref}
\hypersetup{
  colorlinks=true,
  linkcolor=black,
  urlcolor=blue!60!black,
  citecolor=black,
  pdfborder={0 0 0},
  breaklinks=true
}
% Bedre URL-bryting i referanselista (deler URL-er hvor som helst)
\usepackage{xurl}

% --- Avsnittavstand ---
\setlength{\parskip}{0.5em}
\setlength{\parindent}{0pt}

% --- Header og enumitem-fix ---
\providecommand{\tightlist}{\setlength{\itemsep}{2pt}\setlength{\parskip}{2pt}}

% --- Kode-blokker ---
\usepackage{fvextra}
\DefineVerbatimEnvironment{Highlighting}{Verbatim}{breaklines,commandchars=\\\{\}}

% --- Sideskift etter tittel, abstract og ToC ---
\usepackage{etoolbox}
\AtBeginDocument{\pagenumbering{roman}}
\AtEndEnvironment{abstract}{\newpage}
\pretocmd{\tableofcontents}{}{\typeout{ToC}}{}
\apptocmd{\tableofcontents}{\clearpage\pagenumbering{arabic}\setcounter{page}{1}}{\typeout{ToC}}{}

% Tittelside-formattering: litt mer luft og sentrert
\usepackage{titling}
\pretitle{\begin{center}\LARGE\bfseries\color{black!80}}
\posttitle{\par\end{center}\vspace{0.5em}}
\preauthor{\begin{center}\normalsize}
\postauthor{\par\end{center}}
\predate{\begin{center}\normalsize\itshape}
\postdate{\par\end{center}\vspace{2em}}
